\part{Advanced Tools and Administration}

% =====================================================================
\chapter{Multi-paper Comparison}
\label{ch:compare}

After selecting 3--5 papers on the same topic, AI can generate a structured comparison
matrix to help you quickly sort out differences in methods, datasets, metrics, and conclusions.

\section{How to Use}

Select papers in the library (at least 2), then click ``Compare'' in the toolbar. LitFolio
will call the LLM to perform a structured analysis of each paper's TL;DR and deep-read content.

\litfigure{30-compare-en.png}{Multi-paper comparison page. The table is organized into six columns: Problem / Method / Dataset / Key Metric / Limitation /
Conclusion. Cells are directly editable.}

\subsection{Comparison Generation Pipeline}

\begin{enumerate}
  \item \textbf{Material Collection}---for each selected paper, query its TL;DR, key findings,
    deep-read four-section content (Problem / Method / Comparison / Limitations), and abstract.
    If a paper has no deep-read results, only TL;DR + abstract are used.
  \item \textbf{Prompt Construction}---concatenate all paper materials in a fixed format,
    separated by \cmd{=== PAPER: <title> ===} for each paper. The system prompt instructs
    the LLM to return a JSON array, where each element corresponds to one paper and contains
    six columns of content.
  \item \textbf{JSON Parsing}---three-level fault-tolerant parsing (consistent with TL;DR).
    If parsing fails, an attempt is made to extract data from the LLM's raw text in table format.
  \item \textbf{Storage}---results are written to the \cmd{paper\_comparisons} table.
    The \cmd{paper\_ids} field is stored as a JSON array, and the \cmd{content} field stores
    the Markdown table text.
\end{enumerate}

\subsection{Comparison Table Structure}

The comparison results are presented in a table with the following columns:

\begin{longtable}{p{3cm} p{9.5cm}}
\toprule
\textbf{Column} & \textbf{Content} \\
\midrule
\endhead
Paper & Paper title \\
Problem & The problem addressed (2--3 sentences) \\
Method & Core method (focus on innovations) \\
Dataset & Datasets used and their scale \\
Key Metric & Primary metric name and value (e.g., \cmd{BLEU 41.2}, \cmd{F1 0.89}) \\
Limitation & Limitations (1--2 sentences) \\
Conclusion & Conclusion (1 sentence) \\
\bottomrule
\end{longtable}

\subsection{Editing and Saving}

\begin{itemize}
  \item The table supports \textbf{inline editing}---click any cell to modify the AI output.
    Edited content is automatically saved to the database.
  \item Comparison results persist across sessions---no need to regenerate on next open.
  \item The ``Regenerate'' button reruns the AI with the latest paper content (overwriting
    existing results).
  \item Comparing 5 papers typically completes within 30 seconds.
\end{itemize}

\begin{tipbox}
Comparison quality depends on the deep-read quality of each paper. If a paper has not been
deep-read, only the TL;DR and abstract will be used, resulting in less information. It is
recommended to run a deep read on all papers before comparing.
\end{tipbox}

% =====================================================================
\chapter{Markdown Export}
\label{ch:markdown-export}

LitFolio's notes and highlights can be exported as structured Markdown files, compatible
with Obsidian, Logseq, and general Markdown workflows.

\section{Export Content}

Each paper generates a single \cmd{.md} file with the filename format:
\cmd{<first\_author>\_<year>\_<first\_word\_of\_title>.md}. Special characters in
filenames are replaced with underscores.

\subsection{File Structure}

\begin{verbatim}
---
title: "Attention Is All You Need"
authors: ["Vaswani, Ashish", "Shazeer, Noam"]
year: 2017
venue: "NeurIPS"
doi: "10.48550/arXiv.1706.03762"
arxiv_id: "1706.03762"
tags: ["transformer", "attention"]
folders: ["ML/Transformer"]
read_status: "read"
bibtex: "@inproceedings{vaswani2017attention, ...}"
custom_project: "thesis-ch3"
custom_priority: "high"
---

## Notes

Free-text note content (from note.md).

## Structured Notes

### Problem
What problem is being solved...

### Method
What method is proposed...

### Key Numbers
Key experimental data...

### Limitations
Limitations and open problems...

### Thoughts
Your own reflections...

## Highlights

### Key Finding
- **[p.3]** The self-attention mechanism connects all
  positions with a constant number of operations.
  > User comment (if any)

### Method
- **[p.5]** Multi-head attention allows the model to
  attend to information from different representation
  subspaces at different positions.

## Terms

- **Transformer**: A sequence-to-sequence model based
  on the self-attention mechanism, abandoning recurrent
  and convolutional structures.
  - Evidence: "We propose a new simple network
    architecture, the Transformer."
  - Cross-paper reuse: 3
\end{verbatim}

\subsection{Frontmatter Fields}

The YAML frontmatter contains all metadata fields for the paper. Custom metadata fields
are written in the format \cmd{custom\_<field\_name>}. Tags and folders are output as
arrays, making it easy to query with Obsidian's Dataview plugin.

\subsection{Highlights Grouped by Label}

Highlights are grouped by annotation color label (see Section 4.6). Each label serves as
a third-level heading, with its highlights sorted by page number. Each highlight includes:

\begin{itemize}
  \item Page number and color label.
  \item Highlight text (bold).
  \item User comment (if any, displayed as a block quote).
\end{itemize}

\subsection{Terms and Bidirectional Links}

Terms are output with each term name wrapped in \cmd{[[term name]]} double brackets,
leveraging Obsidian's wikilink syntax for graph integration. If a note in the same paper
references a term, Obsidian will automatically establish a bidirectional link.

\section{Incremental Export Algorithm}

Incremental export is tracked via the \cmd{papers.last\_exported\_at} field:

\begin{enumerate}
  \item Query \cmd{SELECT id FROM papers WHERE last\_exported\_at IS NULL
    OR updated\_at > last\_exported\_at}.
  \item For each paper, compare \cmd{updated\_at} (last metadata/note/highlight change time)
    with \cmd{last\_exported\_at} (last export time).
  \item Only export papers where \cmd{updated\_at > last\_exported\_at}.
  \item After successful export, update \cmd{last\_exported\_at = current timestamp}.
\end{enumerate}

This ensures that batch exports do not redundantly write unchanged files.

\section{Configuration and Triggers}

Configure in ``Settings $\rightarrow$ Export'':

\begin{itemize}
  \item \textbf{Export Directory}---e.g., \cmd{\textasciitilde/ObsidianVault/LitFolio/}.
  \item \textbf{Incremental Export}---enabled by default. When disabled, all papers are
    exported each time.
  \item \textbf{Auto Export}---when enabled, saving notes or highlights automatically
    triggers an export (5-second debounce to avoid frequent writes).
\end{itemize}

You can also manually click ``Export Now'' to batch-export all papers. A full export of
500 papers typically completes within 10 seconds.

\begin{tipbox}
Obsidian automatically parses YAML frontmatter. You can filter and sort papers in
Obsidian's property view by fields such as \cmd{year}, \cmd{read\_status}, and \cmd{tags}.
Combined with the Dataview plugin, you can create dynamic query views such as
``all unread papers published after 2024.''
\end{tipbox}

% =====================================================================
\chapter{Command Palette}
\label{ch:palette}

The command palette is LitFolio's keyboard-driven quick-action hub. Press \kbd{Ctrl+K}
(macOS: \kbd{Cmd+K}) on any page to open it.

\litfigure{31-command-palette-en.png}{Command palette. The top input field supports fuzzy search, with results grouped into Navigation / Papers / Actions.}

\section{Three Modes}

The command palette indexes three data sources simultaneously, filtering in real time
as you type:

\begin{description}
  \item[Navigation] Page route names (\cmd{library}, \cmd{reader}, \cmd{graph},
    \cmd{settings}, \cmd{browse}, \cmd{feeds}, \cmd{topic}, \cmd{ask},
    \cmd{import}, \cmd{compare}). Selecting one navigates to that page.
  \item[Search] Paper titles and author names. Queries the database for the 20 most
    recent paper titles and performs fuzzy matching against user input. Selecting one
    opens it in the reader.
  \item[Actions] Executable commands such as \cmd{export bibtex}, \cmd{generate terms},
    \cmd{ai discover}, \cmd{extract concepts}, \cmd{find duplicates}.
    Selecting one directly executes the corresponding action.
\end{description}

\subsection{Fuzzy Matching Algorithm}

The command palette uses a subsequence-based fuzzy matching algorithm rather than simple
substring matching. The algorithm flow:

\begin{enumerate}
  \item \textbf{Input Normalization}---convert both user input and all candidates to
    lowercase and strip non-alphanumeric characters.
  \item \textbf{Subsequence Matching}---check whether each character of the user input
    appears in the candidate in order. For example, \cmd{grpah} against \cmd{graph}:
    g->g checkmark, r->r checkmark, p->a cross, skip a, p->p checkmark,
    a->a checkmark, h->h checkmark $\rightarrow$ match successful.
  \item \textbf{Scoring}---after a successful match, calculate a relevance score:
    \begin{itemize}
      \item Consecutive matching characters earn bonus points (\cmd{graph} matching
        \cmd{graph} scores higher than matching \cmd{generate paragraph}).
      \item Earlier match positions earn more bonus points.
      \item An exact match earns additional bonus points.
    \end{itemize}
  \item \textbf{Sorting}---sort by score in descending order and return the top 10 results.
\end{enumerate}

The advantage of this algorithm is its strong fault tolerance---even if the user makes
a typo (\cmd{librray} $\rightarrow$ \cmd{library}), as long as the characters appear
in order, the match succeeds.

\subsection{Result Grouping and Display}

Results are displayed grouped by category:

\begin{itemize}
  \item \textbf{Recent}---the 5 most recently used commands and 3 most recently opened
    papers, displayed at the top by default (no input required).
  \item \textbf{Navigation}---matched page routes.
  \item \textbf{Papers}---matched papers.
  \item \textbf{Actions}---matched commands.
\end{itemize}

Use the up and down arrow keys to move the selection, \kbd{Enter} to execute, and
\kbd{Esc} to close.

\begin{tipbox}
The command palette is designed so that all features are reachable via keyboard. If you
are unsure where a feature is, searching in the command palette is usually faster than
browsing menus. It also works well as a quick paper-jumping tool---enter a few keywords
from the title to find the target paper.
\end{tipbox}

% =====================================================================
\chapter{Custom Metadata Fields}
\label{ch:custom-fields}

Different researchers have different organizational needs. Custom metadata fields let
you add arbitrary key-value pairs to papers.

\section{Field Management}

Define fields in ``Settings $\rightarrow$ Custom Fields.'' Field definitions are stored
in the \cmd{custom\_field\_defs} table, and field values are stored in the
\cmd{paper\_custom\_fields} table.

\subsection{Field Types}

\begin{longtable}{p{2.5cm} p{4cm} p{6cm}}
\toprule
\textbf{Type} & \textbf{Storage Format} & \textbf{UI Control} \\
\midrule
\endhead
\cmd{text} & Arbitrary string & Text input field \\
\cmd{number} & Numeric string & Number input field, supports sorting comparisons \\
\cmd{date} & Unix timestamp & Date picker \\
\cmd{select} & Option value string & Dropdown select; options defined as a JSON array
  in the \cmd{options} field \\
\bottomrule
\end{longtable}

\subsection{Data Model}

\begin{itemize}
  \item \cmd{custom\_field\_defs}---field definition table. One row per field:
    \cmd{name} (unique), \cmd{field\_type}, \cmd{options} (JSON array of selectable values
    for select type), \cmd{created\_at}.
  \item \cmd{paper\_custom\_fields}---field value table. Composite primary key
    \cmd{(paper\_id, field\_id)}; each row stores one field value for one paper.
\end{itemize}

\section{Usage}

Once fields are defined, a ``Custom Fields'' section appears in every paper's detail
drawer. You can directly edit each paper's field values. Custom fields support:

\begin{itemize}
  \item \textbf{Filtering}---in the main list's filter bar, you can filter papers by
    custom field values. For example, filter all papers where
    \cmd{project = thesis-ch3}.
  \item \textbf{Sorting}---number and date type fields support ascending/descending sort.
  \item \textbf{Export}---in Markdown export's YAML frontmatter, custom fields are
    written in the format \cmd{custom\_<field\_name>}.
  \item \textbf{Global Definitions}---define once, and all papers share the same field
    definitions. Deleting a field definition cascades to delete that field's values
    across all papers.
\end{itemize}

\begin{tipbox}
The difference between custom fields and tags: tags are flat, untyped keywords; custom
fields are typed, structured data (such as dates, numbers, enums). If you need to sort
by ``deadline'' or filter by ``priority,'' custom fields are more appropriate than tags.
\end{tipbox}

% =====================================================================
\chapter{Settings}
\label{ch:settings}

The \cmd{/settings} route is LitFolio's ``control panel.'' The main sections are:
\textbf{LLM Configuration}, \textbf{Task Assignment}, \textbf{WebDAV Sync},
\textbf{Export}, \textbf{Custom Fields}, \textbf{Annotation Colors}, and
\textbf{Topic Alerts}.

\section{LLM Configuration}

LitFolio is not locked to any single LLM provider. Its design follows the ``OpenAI-compatible
protocol'': you can add any number of endpoint configurations (profiles), with one marked
as ``active.'' Each profile contains an API URL, API key, chat model, embedding model
(optional), and sampling temperature.

\litfigure{17-settings-profiles-en.png}{LLM Configuration. Two profiles: DeepSeek and local Ollama; the currently active profile is DeepSeek.
Each profile can ``Pull'' the remote available model list to populate the chatModel field.}

\subsection{Pulling Available Models}

After filling in the API URL and key, click \cmd{Pull}. LitFolio will call the remote
\cmd{GET /v1/models} endpoint to fetch the available model list. This step prevents
manual typos in model names. The returned models populate the dropdown.

\subsection{Testing the Connection}

``Test Connection'' initiates a small chat request (prompt: ``Hello'') and verifies the
entire pipeline using a real model response. Failure reasons (authentication, network,
model not found) are displayed as-is.

\subsection{Reasoning Model Compatibility}

LitFolio automatically adapts to OpenAI reasoning-series models (GPT-5, o1, o3, o4, etc.):
these models use \cmd{max\_completion\_tokens} instead of \cmd{max\_tokens} and do not
support the \cmd{temperature} parameter. LitFolio switches the request format automatically
based on the model name---no manual configuration needed. If you use one of these models,
simply enter the model name in the profile as usual.

\section{Task Assignment}

LitFolio groups all LLM calls into \textbf{6 tasks}:

\begin{itemize}
  \item \textbf{Quick Read TL;DR}---one-sentence summary + key findings
  \item \textbf{Deep Read quickRead}---four-section deep reading
  \item \textbf{Translation translate}---title/abstract/selection translation
  \item \textbf{Search Expansion search}---query expansion for topic search
  \item \textbf{Survey survey}---topic survey generation
  \item \textbf{Terms term}---glossary extraction
  \item \textbf{Ask ask}---in-library RAG answer
\end{itemize}

Each task can be independently bound to any profile and any specific model. This lets
you configure:
\begin{description}
  \item[Expensive but Smart] Deep read, survey, ask \ensuremath{\rightarrow} choose
    \textbf{top-tier proprietary chat models} (such as the latest flagships from OpenAI,
    Anthropic, Google). These affect answer quality, but per-call token usage is modest.
  \item[cheap and Capable] Translation, terms \ensuremath{\rightarrow} choose
    \textbf{budget tiers from major domestic API providers} (DeepSeek, Doubao, Zhipu,
    Moonshot, etc., all offer excellent value). Translation/terms have high call volume,
    but at a few yuan per million tokens, costs stay minimal.
  \item[Local and Free] TL;DR, search expansion \ensuremath{\rightarrow} \textbf{open-source
    chat models on local Ollama} (Qwen, Llama, DeepSeek distilled versions, 7B--32B
    all work). Zero cost, good privacy, speed depends on your GPU.
\end{description}

\subsection{LLM Call Details per Task}

Understanding how each task calls the LLM helps you decide what tier of model to assign.

\subsubsection{Quick Read TL;DR}

Input: title + authors (first 6) + abstract. The prompt requests JSON: \cmd{\detokenize{{"tldr": "one
sentence, no more than 40 words", "key\_findings": ["3--5 items, each no more than 20
words"]}}}. Output language follows user settings. JSON parsing has three-level fault
tolerance (direct parse $\rightarrow$ extract \cmd{\{...\}} substring $\rightarrow$ empty
object fallback). Results are written to the \cmd{papers} table cache; clicking again
will not re-invoke the LLM.

\subsubsection{Deep Read Quick Read}

Input is the same as TL;DR (title + authors + abstract), but the prompt is more detailed:
it requests four-section structured JSON (problem / method / comparison / limitations),
each section 2--4 sentences in prose style, no bullet points allowed. The system prompt
positions the LLM as ``a senior reviewer helping a colleague decide whether to read the
full paper,'' guiding it to provide judgmental assessments rather than generic summaries.

\subsubsection{Translation}

Input: title + abstract. The prompt instructs the LLM to preserve technical terms in
their original language (model names, algorithm names, mathematical symbols, units,
dataset names), not to paraphrase or summarize, and return JSON
\cmd{\detokenize{{"title": "...", "abstract": "..."}}}.

\section{WebDAV Sync}

\litfigure{18-settings-sync-en.png}{WebDAV Sync panel. After configuring the URL and credentials, you can \textbf{Push} (local$\rightarrow$remote) or \textbf{Pull and Replace Local}
(remote$\rightarrow$local, which first backs up the local state to \fpath{backups/sync/}).}

LitFolio's sync uses ``full library snapshot'' semantics, not fine-grained incremental
sync. Specifically:

\begin{enumerate}
  \item \textbf{Push}---packages the entire \fpath{\textasciitilde/Litera-Library/} and
    pushes it to WebDAV, creating a complete copy on the remote.
  \item \textbf{Pull and Replace Local}---pulls a complete copy from the remote and
    overwrites the local state.
\end{enumerate}

\begin{warnbox}
``Pull and Replace Local'' will \textbf{overwrite} the local database and all
\fpath{papers/}. LitFolio automatically saves the current local state to
\fpath{backups/sync/<timestamp>/} before executing, so you can roll back even if you
pull in the wrong direction---but please take this operation seriously, especially when
both machines have unpushed changes.
\end{warnbox}

The design choice of ``full library overwrite'' rather than incremental merging is because
most modifications to a paper library involve PDFs and notes, and the semantics of which
side wins during a conflict are difficult to explain to users; exposing only ``push all /
pull all'' with ``automatic backup'' is more predictable.

\section{Export Settings}

The ``Settings $\rightarrow$ Export'' panel controls Markdown export behavior (see
Chapter 14 for details):

\begin{itemize}
  \item \textbf{Export Directory}---choose the output location for \cmd{.md} files.
  \item \textbf{Incremental Export}---enabled by default. When disabled, all papers are
    exported each time.
  \item \textbf{Auto Export}---when enabled, saving notes or highlights automatically
    triggers an export (5-second debounce).
  \item \textbf{Export Now}---manually trigger a full or incremental export.
\end{itemize}

\section{Annotation Color Settings}

The ``Settings $\rightarrow$ Annotation Colors'' panel manages the mapping between
highlight colors and semantic labels (see Section 4.6 for details):

\begin{itemize}
  \item View all current color-label mappings.
  \item Edit label names.
  \item Add new color-label mappings.
  \item Delete unused mappings.
\end{itemize}

\section{Custom Field Settings}

The ``Settings $\rightarrow$ Custom Fields'' panel manages custom metadata field
definitions for papers (see Chapter 16 for details):

\begin{itemize}
  \item Create new fields (name + type + optional select options).
  \item Edit existing fields.
  \item Delete fields (cascading deletion of that field's values across all papers).
\end{itemize}

\section{Topic Alert Settings}

The ``Settings $\rightarrow$ Topic Alerts'' panel manages automated paper monitoring
(see Section 7.4 for details):

\begin{itemize}
  \item View all alerts and their last run times.
  \item View the history of results for each alert.
  \item Edit an alert's query, frequency, and target folder.
  \item Pause or delete alerts.
  \item Mark results as read.
\end{itemize}

% =====================================================================
\chapter{Data and File Layout}
\label{ch:data}

LitFolio believes ``your data belongs to you.'' All information is stored in readable
formats under \fpath{\textasciitilde/Litera-Library/}, so you can access it with scripts
even without LitFolio.

\section{Top-Level Structure}

\begin{quote}\small\ttfamily
\textasciitilde/Litera-Library/\\
\ \ library.db\hfill\textnormal{\small Main SQLite database}\\
\ \ library.db-wal\hfill\textnormal{\small WAL log (runtime)}\\
\ \ library.db-shm\hfill\textnormal{\small Shared memory (runtime)}\\
\ \ papers/\\
\ \ \ \ \textless paperId\textgreater /\\
\ \ \ \ \ \ paper.pdf\\
\ \ \ \ \ \ meta.json\\
\ \ \ \ \ \ note.md\\
\ \ \ \ \ \ pdf-text.txt\\
\ \ notes/\\
\ \ \ \ ask/\hfill\textnormal{\small Ask results saved as notes}\\
\ \ backups/\\
\ \ \ \ deleted/\hfill\textnormal{\small Paper copies before deletion}\\
\ \ \ \ sync/\hfill\textnormal{\small Local snapshots before pull}\\
\ \ sync/\hfill\textnormal{\small Temporary area for WebDAV push}
\end{quote}

\section{library.db Table Overview}

Main tables (detailed schema in the repository at \fpath{src-tauri/migrations/}, totaling
24 migration files):

\subsection{Core}

\begin{itemize}
  \item \cmd{papers}---core metadata, one row per paper (id, title, authors, abstract, status,
    tldr, quick\_read\_json, translated\_title, translated\_abstract, bibtex,
    last\_exported\_at, ...)
  \item \cmd{highlights}---selection highlights and area highlights (paper\_id, page,
    boundingRect, text, comment, color, label)
  \item \cmd{terms}---glossary (paper\_id, term, definition, evidence, cross\_refs)
  \item \cmd{tags} / \cmd{paper\_tags}---tags and associations
  \item \cmd{folders} / \cmd{paper\_folders}---folders and associations (supports nesting
    and multiple membership)
\end{itemize}

\subsection{Discovery and Reading}

\begin{itemize}
  \item \cmd{feeds} / \cmd{feed\_items}---RSS feeds and entries
  \item \cmd{topic\_surveys}---saved topic surveys
  \item \cmd{topic\_alerts} / \cmd{topic\_alert\_results}---topic alert configurations
    and results
  \item \cmd{reading\_queue}---reading queue (paper\_id, priority, target\_date, note)
  \item \cmd{paper\_comparisons}---multi-paper comparison results
  \item \cmd{recommendation\_cache}---similar paper recommendation cache
  \item \cmd{paper\_citations}---citation network cache
\end{itemize}

\subsection{Notes and Knowledge}

\begin{itemize}
  \item \cmd{paper\_note\_sections}---structured note cards (paper\_id, section\_key,
    content, source, sort\_order)
  \item \cmd{paper\_links}---manual/AI relationships between papers
  \item \cmd{paper\_terms}---paper-term associations (including cross-paper reuse)
  \item \cmd{concepts}---concept nodes (name, description, source)
  \item \cmd{concept\_relations}---relationships between concepts
    (replaces/extends/requires/enables/competes\_with)
  \item \cmd{paper\_concepts}---paper-concept associations
\end{itemize}

\subsection{Configuration and Indexing}

\begin{itemize}
  \item \cmd{llm\_profiles} / \cmd{task\_assignments}---endpoints and task bindings
  \item \cmd{papers\_fts} / \cmd{highlights\_fts} / \cmd{terms\_fts}---FTS5 full-text indexes
  \item \cmd{paper\_embeddings}---vector embedding cache
  \item \cmd{custom\_field\_defs} / \cmd{paper\_custom\_fields}---custom metadata
  \item \cmd{smart\_collections}---smart collection rules
\end{itemize}

\section{papers/<id>/ Per-Paper Layout}

Each paper has its own directory, named with a ULID (a 26-character string like
\cmd{01KSCN65XF4B2PZ83D27ET55PX}):

\begin{description}
  \item[paper.pdf] The original PDF. If this file is deleted, LitFolio will display
    ``PDF unbound,'' but metadata and notes remain intact; use ``Rebind PDF'' in the
    library drawer to reattach it.
  \item[meta.json] Metadata snapshot (title, authors, year, DOI/arXiv ID, source, etc.).
    SQLite is the primary storage; meta.json is a mirror for script access.
  \item[note.md] Notes for this paper. Entirely user-written; LitFolio does not modify
    content beyond what you write in the reader.
  \item[pdf-text.txt] Full-text extraction from the PDF (generated offline with
    pdf-extract), used for FTS5 indexing and RAG retrieval.
\end{description}

\section{Performance Optimization}

LitFolio includes two key optimizations for large-scale libraries (1000+ papers).

\subsection{Virtual Scrolling}

The library main list uses virtual scrolling, implemented with
\cmd{@tanstack/react-virtual}. Core principles:

\begin{enumerate}
  \item \textbf{Container Measurement}---when the list container mounts, measure the
    visible area height (e.g., 800px).
  \item \textbf{Row Height Estimation}---estimate the height of each paper card row
    (e.g., 72px).
  \item \textbf{Viewport Calculation}---based on scroll offset and container height,
    calculate the range of row indices that need to be rendered. For example, when
    scrolled to 500px, render rows 7--18 (plus an overscan buffer of 3 rows above
    and below).
  \item \textbf{DOM Node Pool}---only create DOM nodes within the visible range. As the
    user scrolls, nodes leaving the viewport are recycled and new nodes entering the
    viewport are created.
  \item \textbf{Placeholder Padding}---the total list height is calculated as
    \cmd{totalSize $\times$ rowHeight}, using top and bottom padding elements to
    maintain the scrollbar, ensuring scroll behavior is consistent with a full list.
\end{enumerate}

Results:

\begin{itemize}
  \item Initial render time $< 100$ms (2000 papers).
  \item Scrolling maintains 60fps---because only 20--30 DOM nodes are manipulated at
    a time.
  \item Memory usage is proportional to the visible area size, not to the total number
    of papers.
\end{itemize}

\subsection{Embedding Cache}

RAG queries (in-library asking) require vector embeddings of paper text. LitFolio uses
the \cmd{paper\_embeddings} table to cache generated embedding vectors:

\begin{enumerate}
  \item Before a query, check whether the \cmd{paper\_embeddings} table has a cached
    record for the paper and whether the \cmd{content\_hash} matches the current text hash.
  \item Cache hit $\rightarrow$ use the cached embedding vector directly, skipping the
    API call.
  \item Cache miss $\rightarrow$ call the embedding API to generate the vector and write
    it to the cache table.
  \item Paper content changes (note edits, highlight additions/deletions) $\rightarrow$
    \cmd{content\_hash} changes $\rightarrow$ embedding is automatically regenerated
    on the next query.
\end{enumerate}

The embedding cache is fully transparent to users---no manual operation is required.
Logs will show \cmd{"embedding cache hit"} or \cmd{"embedding cache miss"} to confirm
whether the cache is working.

\section{Backups}

LitFolio has two types of automatic backups:

\begin{enumerate}
  \item \textbf{Deletion Recycle Bin} (\fpath{backups/deleted/})---as described in
    \textbf{Chapter 2}, deleted papers are moved here first. Automatically cleaned
    after 30 days.
  \item \textbf{Sync Snapshots} (\fpath{backups/sync/})---a complete local snapshot
    taken before ``Pull and Replace.'' Retains the most recent 5 snapshots.
\end{enumerate}

If you need to add your own daily scheduled backups, the simplest approach is to use
the entire \fpath{\textasciitilde/Litera-Library/} as a borg/restic repository
source---the SQLite WAL files are included, ensuring consistency.

% =====================================================================
\chapter{Keyboard Shortcuts Quick Reference}
\label{ch:shortcuts}

\section{Reader (\cmd{/reader})}

\begin{longtable}{p{5.5cm} p{8cm}}
\toprule
\textbf{Shortcut} & \textbf{Action} \\
\midrule
\endhead
\kbd{Ctrl+F} / \kbd{Cmd+F} & Open PDF search box \\
\kbd{Enter} & Jump to next match \\
\kbd{Shift+Enter} & Jump to previous match \\
\kbd{Esc} & Close search box \\
\kbd{Alt+Drag} & Area highlight (rectangular screenshot) \\
\kbd{Mouse Selection} & Text highlight (bubble appears) \\
\kbd{j} / \kbd{k} & Next page / Previous page \\
\kbd{1} / \kbd{2} / \kbd{3} & Switch right panel: Notes / Selection Translation / Terms \\
\kbd{[} / \kbd{]} & Left/Right panel width 5\% step \\
\bottomrule
\end{longtable}

\section{Global}

\begin{longtable}{p{5.5cm} p{8cm}}
\toprule
\textbf{Action} & \textbf{Effect} \\
\midrule
\endhead
\kbd{Ctrl+K} / \kbd{Cmd+K} & Open command palette \\
Drag PDF to any page & Navigate to ``Import $\rightarrow$ PDF File'' with the file prefilled \\
\kbd{Ctrl+Enter} / \kbd{Cmd+Enter} & Send message on Ask page \\
\kbd{Esc} & Close command palette / dialog / search box \\
\bottomrule
\end{longtable}

% =====================================================================
\chapter{Troubleshooting}
\label{ch:troubleshoot}

Below are the most commonly encountered issues, organized by ``Symptom $\rightarrow$
Cause $\rightarrow$ Fix.''

\section{LLM Call Failure}

\textbf{Symptom}: TL;DR / Deep Read / Translation buttons spin and then show
``Request failed.''

\textbf{Common Causes and Fixes}:

\begin{enumerate}
  \item The current profile is not configured or the model name is misspelled---go to
    ``Settings $\rightarrow$ LLM Configuration'' and click ``Test Connection.''
  \item API key balance exhausted / authentication expired---error details are displayed
    as-is; refer to your provider's instructions.
  \item Network unreachable (especially direct access to OpenAI from certain regions)---switch
    to a compatible endpoint (Azure / regional proxy).
  \item Local Ollama is not running---start with \cmd{ollama serve}, then go to settings
    and test.
\end{enumerate}

\section{PDF Load Failure}

\textbf{Symptom}: The reader shows a perpetual spinner or displays ``Failed to load PDF.''

\textbf{Fix}:

\begin{itemize}
  \item Check whether \fpath{papers/<id>/paper.pdf} exists; if not, click ``Rebind PDF''
    in the library drawer.
  \item Some scanned PDFs use non-standard encodings that cause PDF.js rendering to
    fail---try processing with \cmd{ocrmypdf} or \cmd{qpdf --linearize} and then rebind.
\end{itemize}

\section{RSS Not Fetching / Always ``Already Up to Date''}

\textbf{Symptom}: The remote clearly has new content, but clicking ``Refresh'' shows
``Already up to date.''

\textbf{Cause}: LitFolio uses conditional GET (\cmd{If-Modified-Since} / \cmd{ETag});
if the remote says nothing has changed, it will not re-fetch. If you suspect the remote
is reporting incorrectly, delete the feed from ``Feed Management'' and re-add it.

\textbf{User-Agent Blocked}: A small number of journal websites block the default UA.
Enter a custom \cmd{User-Agent} header when subscribing.

\section{WebDAV Sync Failure}

Three common scenarios:

\begin{itemize}
  \item \textbf{401 Unauthorized}---incorrect username or password. For Nextcloud, use
    an \textbf{app password} rather than your account password.
  \item \textbf{404 / 405}---incorrect WebDAV URL path. Nextcloud uses
    \cmd{/remote.php/dav/files/<user>/}, and Jianguo Cloud uses
    \cmd{https://dav.jianguoyun.com/dav/}.
  \item \textbf{Conflict / Midway Failure}---the local snapshot taken before the pull
    is in \fpath{backups/sync/} and can be manually restored.
\end{itemize}

\section{Database Locked}

\textbf{Symptom}: On startup, a ``database is locked'' error appears.

\textbf{Cause}: After an unexpected kill, WAL files remain. In most cases, SQLite will
automatically checkpoint and recover. If LitFolio fails to start repeatedly, close all
LitFolio processes and delete \fpath{library.db-wal} and \fpath{library.db-shm} (the
data itself is in \fpath{library.db}; deleting the WAL only loses the most recent few
seconds of unpersisted changes).

\section{Partial Batch Import Failure}

\textbf{Symptom}: When importing a folder, some PDFs show ``Failed.''

\textbf{Common Causes}:

\begin{enumerate}
  \item \textbf{Corrupt PDF}---the file is incomplete or truncated. Verify with
    \cmd{qpdf --check}.
  \item \textbf{No Text Layer}---scanned PDF without an OCR text layer. LitFolio will
    still import it (using the filename as the title), but full-text search will not
    cover its content. Process with \cmd{ocrmypdf} and then rebind the PDF.
  \item \textbf{Permission Issues}---the PDF file is locked by another program or
    lacks read permissions.
\end{enumerate}

\section{Linux AppImage Won't Launch}

\textbf{Symptom}: Double-clicking the \fpath{.AppImage} file does nothing, or running
from the terminal produces an error.

\textbf{Common Causes and Fixes}:

\begin{enumerate}
  \item \textbf{Missing execute permission}---run
    \cmd{chmod +x LitFolio\_x.y.z\_amd64.AppImage}.
  \item \textbf{FUSE not installed}---error message contains ``fuse'' or ``libfuse''.
    Ubuntu/Debian: \cmd{sudo apt install libfuse2};
    Fedora: \cmd{sudo dnf install fuse-libs}.
  \item \textbf{Sandbox conflict}---on some desktop environments (certain Wayland
    compositors), the Chromium sandbox may fail. Try launching with the
    \cmd{-{}-no-sandbox} flag from the terminal.
  \item \textbf{Insufficient disk space}---AppImage extracts to \fpath{/tmp} at
    runtime; ensure at least 500\,MB of free space.
\end{enumerate}

\section{Concept Extraction Returns Empty}

\textbf{Symptom}: After clicking ``Extract Concepts,'' no concept nodes appear in the
knowledge graph.

\textbf{Common Causes}:

\begin{enumerate}
  \item The paper has no TL;DR or deep-read results---the LLM lacks sufficient context
    to identify concepts. Run a quick read or deep read on the paper first.
  \item The LLM returned unparseable JSON---check whether the model configured in your
    LLM settings supports JSON output format.
  \item The paper content is too short or is not an academic paper---concept extraction
    works best on academic papers.
\end{enumerate}

\section{Citation Network / Similar Papers Empty}

\textbf{Symptom}: After clicking ``Citations'' or ``Similar Papers,'' it shows ``No
data available.''

\textbf{Cause}:

\begin{itemize}
  \item The paper has no DOI and its title cannot be matched to a record in Semantic
    Scholar.
  \item The paper is very new (published less than 1 week ago) and has not yet been
    indexed by S2.
  \item Niche field or non-English paper---S2 coverage is limited in some areas.
  \item Network issues---check whether \cmd{api.semanticscholar.org} is accessible.
\end{itemize}

\section{Smart Collection Not Updating}

\textbf{Symptom}: After changing a paper's tags or status, the smart collection's
paper list does not change.

\textbf{Cause}: Smart collections are queried in real time---you need to re-select the
collection to see the update. This is not a bug but a design choice: avoiding rule
queries triggered by every paper change to prevent performance impacts.

% =====================================================================
\appendix

% ---------------------------------------------------------------------
\chapter{LLM Endpoint Compatibility List}
\label{appx:llm}

LitFolio follows the ``OpenAI Chat Completions protocol.'' The following endpoints have
been verified to work directly:

\begin{longtable}{p{4cm} p{5.5cm} p{4.5cm}}
\toprule
\textbf{Service} & \textbf{API URL} & \textbf{Notes} \\
\midrule
\endhead
OpenAI Official & \cmd{https://api.openai.com/v1} & Direct access may be restricted \\
DeepSeek & \cmd{https://api.deepseek.com/v1} & Recommended: great value \\
Azure OpenAI & \cmd{<resource>.openai.azure.com/} \mbox{}\cmd{openai/deployments/<dep>/v1} & Requires API version header \\
Ollama (local) & \cmd{http://localhost:11434/v1} & Fully offline \\
LM Studio (local) & \cmd{http://localhost:1234/v1} & Desktop GUI loads GGUF \\
OpenRouter & \cmd{https://openrouter.ai/api/v1} & One key, multiple models \\
Together.ai & \cmd{https://api.together.xyz/v1} & Open-source model hosting \\
\bottomrule
\end{longtable}

\textbf{Task Assignment Recommendations} (specific models change rapidly; follow the
approach and pick what is currently best):

\begin{itemize}
  \item \textbf{Student / Fully Free}---run everything possible locally (local Ollama
    with 7B--14B open-source models is sufficient for TL;DR / translation / search
    expansion). Reserve online calls for ``deep read,'' ``survey,'' and ``ask,''
    using free tiers from domestic providers.
  \item \textbf{Researcher / Best Value}---point all tasks to the \textbf{budget tier
    of a major domestic API} (DeepSeek, Doubao, Zhipu, Moonshot---pick one). Low unit
    price, large context window, sufficient for the vast majority of research scenarios.
  \item \textbf{Power User / No Compromise}---assign TL;DR / translation / terms to
    a domestic budget tier; assign deep read / survey / ask to the \textbf{strongest
    flagship chat model available at the time} (OpenAI, Anthropic, Google take turns
    leading---pick one based on release date) to guarantee the highest quality answers.
\end{itemize}

% ---------------------------------------------------------------------
\chapter{arXiv Main Category Quick Reference}
\label{appx:arxiv}

\begin{description}
  \item[cs.AI / cs.LG / cs.CL] Artificial Intelligence / Machine Learning / Natural
    Language Processing
  \item[cs.CV] Computer Vision
  \item[cs.CR / cs.DC] Cryptography and Security / Distributed Systems
  \item[stat.ML] Statistical Machine Learning
  \item[physics.optics] Optics and Photonics
  \item[physics.atom-ph] Atomic Physics
  \item[cond-mat.mes-hall] Condensed Matter -- Mesoscopic and Nanoscale
  \item[quant-ph] Quantum Physics
  \item[math.OC] Optimization and Control
  \item[math.AP] Partial Differential Equations
  \item[q-bio.NC] Neuroscience
\end{description}

For the complete list, see \cmd{https://arxiv.org/category\_taxonomy}.

% ---------------------------------------------------------------------
\chapter{Recommended RSS Feeds}
\label{appx:rss}

\textbf{arXiv Subcategories} (replace \cmd{CAT} with e.g. \cmd{cs.LG}):

\begin{itemize}
  \item \cmd{http://export.arxiv.org/rss/CAT}---official RSS, refreshed daily
  \item \cmd{http://export.arxiv.org/rss/CAT.recent}---latest only
\end{itemize}

\textbf{Journals}:

\begin{itemize}
  \item Nature: \cmd{https://www.nature.com/nature.rss}
  \item Science: \cmd{https://www.science.org/rss/news\_current.xml}
  \item Nature Photonics: \cmd{https://www.nature.com/nphoton.rss}
  \item Optica (OSA): \cmd{https://www.optica-opn.org/rss/news.xml}
  \item PRL: \cmd{http://feeds.aps.org/rss/recent/prl.xml}
\end{itemize}

\textbf{Research Blogs / Surveys}:

\begin{itemize}
  \item Distill: \cmd{https://distill.pub/rss.xml} (no longer updated, but archived
    articles remain valuable)
  \item Lil'Log: \cmd{https://lilianweng.github.io/index.xml}
\end{itemize}

% ---------------------------------------------------------------------
\chapter{License}
\label{appx:license}

LitFolio is released under the \textbf{GPL-3.0-or-later} license. This means:

\begin{itemize}
  \item You are free to use, modify, and redistribute.
  \item If you create and distribute a derivative work based on LitFolio, the derivative
    must also be released under the same license (or a compatible more permissive license)
    and include source code.
  \item No warranty of any kind, express or implied, is provided.
\end{itemize}

The full license text is in the \fpath{LICENSE} file in the repository root directory.
